\documentclass[book.tex]{subfiles}
\begin{document}

\chapter{Lagrangian Neural Networks}
\label{chap:lagrangian_nn}
\noindent\rule{11cm}{0.4pt} \\
\textbf{Prerequisites:} \autoref{chap:lagrangian_mechanics}, \autoref{chap:hamiltonian_nn} \\
\textbf{Difficulty Level:} ** \\
\noindent\rule{11cm}{0.4pt}
\vspace{5mm}

One of the major limitations of Hamiltonian neural networks is that we need to know the canonical coordinates of the system. For many complex systems, this assumption is unrealistic. The Lagrangian neural network framework provides a method to learn a Lagrangian for a system directly from data. This method has the powerful advantage that it can be applied to systems with unknown conservation laws.

\section{Mathematical Framework}

Suppose that we have a physical system that evolves from $x_0 = (q_0, \dot{q_0})$ to $x_1=(q_1, \dot{q_1})$. The system can take many different paths from $x_0$ to $x_1$. The Lagrangian formalism assigns to each such path a score by using the action functional
\begin{align}
    S = \int_{t_0}^{t_1} (T(q_t, \dot{q_t}) - V(q_t)) dt
\end{align}
where $T$ is the kinetic energy and $V$ is the potential energy. The quantity $T - V$ is defined to be the Lagrangian.
\begin{align}
    \mathcal{L} = T - V
\end{align}
The Lagrangian formalism tells us that the path which is chosen will have minimal energy and have the functional derivative equal 0.
\begin{align}
    \partial S = 0
\end{align}
As we saw earlier, the path which satisfies this constraint satisfies the Euler-Lagrange equations
\begin{align}
    \frac{d}{dt}\frac{\partial \mathcal{L}}{\partial \dot{q_j}} &= \frac{\partial \mathcal{L}}{\partial q_j}
\end{align}
Classically, physicists write down an analytical expression for $\mathcal{L}$ (as we have done in previous chapters) and use mathematical techniques to solve for the resulting differential equations. However, for complex systems, we may not know the form of $\mathcal{L}$ but may instead have just raw data readouts from the system. The learning task we face here is to learn $\mathcal{L}$ from data.

\section{Parameteric Lagrangians}

We start by writing down the Euler-Lagrange equations in vectorized form

\begin{align}
    \frac{d}{dt} \nabla_{\dot{q}} \mathcal{L} &= \nabla_q \mathcal{L}
\end{align}
Here we define
\begin{align}
    (\nabla_{\dot{q}} \mathcal{L})_i &= \frac{\partial \mathcal{L}}{\partial \dot{q}_i}
\end{align}
We can apply the chain rule to expand out the derivative %(\textcolor{red}{TODO: This derivation doesn't make full sense. Maybe expand})
\begin{align}
\frac{d}{dt} \left ( \nabla_{\dot q} \mathcal{L}(q, \dot q, t) \right ) &= \nabla_q \mathcal{L} \\
\nabla_{\dot q} \left (\frac{d}{dt} \mathcal{L}(q, \dot q, t) \right ) &= \nabla_q \mathcal{L} \\
\nabla_{\dot q} \left ( \frac{\partial}{\partial q} \mathcal{L}(q, \dot q, t) \dot q + \frac{\partial}{\partial \dot q} \mathcal{L}(q, \dot q, t) \Ddot{q}  \right ) &= \nabla_q \mathcal{L} \\
(\nabla_{\dot{q}} \nabla_{\dot{q}}^T \mathcal{L}) \ddot{q} + (\nabla_q \nabla_{\dot{q}}^T \mathcal{L}) \dot{q}  &= \nabla_q \mathcal{L}
\end{align}
We can apply a matrix inversion to recover the acceleration $\ddot{q}$
\begin{align}
    \ddot{q} &= (\nabla_{\dot{q}} \nabla_{\dot{q}}^T \mathcal{L})) ^{-1} [\nabla_q \mathcal{L} - (\nabla_q \nabla_{\dot{q}}^T \mathcal{L})\dot{q}  ]
\end{align}

%\textcolor{red}{TODO: Find a good derivation of the loss function and add here.}
Assume that we model the Lagrangian $L$ by a neural network
\begin{align}
L &= f_\theta
\end{align}
The gradient operations can be implemented by autograd, so we can directly implement the formula above to compute $\ddot{q}$. To recover $q$ from $\ddot{q}$, we can integrate
\begin{align}
\hat{q} &= \mathrm{Integrate}(\mathrm{Integrate}(\ddot{q}))
\end{align}
The loss function then directly penalizes the difference between the original observations $q$ and the predictions $\hat{q}$.
\begin{align}
\mathcal{L} &= \sum_i \|q_i - \hat{q}_i\|^2
\end{align}

\section{A Particle Falling Under Gravity}

We start with a simple example of a particle falling in gravity
\begin{align}
    L &= \frac{1}{2} m(\dot{q_1}^2 + \dot{q_2}^2 ) - m ggq_1
\end{align}

We then note that
\begin{align}
\nabla_q L &= \begin{pmatrix}
    -mg \\
    0 
    \end{pmatrix} \\
\nabla_{\dot{q}} L &= \begin{pmatrix}
    m \dot{q_1} \\
    m \dot{q_2}
    \end{pmatrix}
\end{align}
Applying these identities, we can compute
\begin{align}
\nabla_{\dot{q}} \nabla_{\dot{q}}^T L &= \begin{pmatrix}
    m & 0 \\
    0 & m 
    \end{pmatrix} \\
\nabla_q \nabla_{\dot{q}}^T &= 0 \\
\nabla_q L &= \begin{pmatrix}
-mg \\
0
\end{pmatrix}
\end{align}

Plugging in values, we then find that
\begin{align}
\begin{pmatrix}
\ddot{q_1} \\
\ddot{q_2}
\end{pmatrix}
&= \begin{pmatrix}
-g \\
0 
\end{pmatrix}
\end{align}

%\textcolor{red}{TODO: Does the derivation check out?}

%\section{Relativistic Motion}

%Let us now consider an example of a relativistic particle with Lagrangian

%\begin{align}
%L &= \left ( \frac{1}{\sqrt{1-\dot{q}^2}} - 1\right)  + gq
%\end{align}


\end{document}